% Section: P5P8-SYNTH D21 Per-Reward-Decile Cross-Pillar Decision Concordance
% (iter 196 JOB B). Fresh 21st density domain (NOT in any prior D1..D20
% SYNTH row). D20 measured aggregate cross-pillar Spearman ρ (4 N2 methods
% × 40 steps = 160 cells) and found a two-cluster structure: {P5, P8}
% cluster (gift wins) ↔ {P6, P7} cluster (areal wins) with NEGATIVE
% cross-cluster ρ. D21 lifts the lens to per-decile granularity.

\subsubsection{Falsifiable questions and operational stakes}\label{sec:synth-d21-questions}
D20 reported an aggregate two-cluster structure: pillars $\{P_5, P_8\}$
(operationally relevant for deployment) cluster with positive
within-cluster $\rho$ and cluster $\{\mathrm{P}_6, \mathrm{P}_7\}$
(registry/controller-relevant for design) cluster with positive
within-cluster $\rho$, with NEGATIVE cross-cluster $\rho$. This finding
was based on 160 cells (4 methods $\times$ 40 steps) on a single N2
same-stack panel. The natural follow-up --- \textbf{is the two-cluster
structure uniform across the reward distribution?} --- has not been
asked. D21 stratifies the 160 cells into 10 reward-deciles based on
$\mathrm{reward\_mean}$ and asks: does the D20 structure break down in
specific deciles? Is there a regime-conditional best method?

\begin{enumerate}
\item Does the D20 two-cluster structure hold within reward-deciles,
  or does it depend on pooling?
\item Does the best method per pillar vary across deciles (regime
  conditional) or is it constant (universal)?
\item Are low-reward deciles (where methods diverge most) more
  informative than high-reward deciles (where methods are saturated)?
\end{enumerate}

\subsubsection{Pipeline and headline findings}\label{sec:synth-d21-pipeline}
We load \texttt{n2\_metrics.tsv} (160 rows: 4 methods $\times$ 40 steps).
Per-(method, step) we compute 4 pillar headliners:
$\mathrm{P}_5 = \overline{\mathrm{reward}}$,
$\mathrm{P}_6 = -\mathrm{ZVF}$,
$\mathrm{P}_7 = \overline{\mathrm{reward}} / (1 + \mathrm{cv\_len})$,
$\mathrm{P}_8 = \overline{\mathrm{reward}} / \overline{\mathrm{len}}$.
We stratify the 160 cells into 10 reward-deciles using the
$\mathrm{reward\_mean}$ quantile. Per-(decile, method) we aggregate by
mean over the steps in the decile, giving 4 method-points per pillar per
decile. We compute point Spearman $\rho$ on these 4 method-points and
bootstrap 95\% CIs ($B{=}2000$) by resampling (method, step) cells within
each decile.

\subsubsection{Headline results}\label{sec:synth-d21-results}
\textbf{1 PASS + 4 sharp FAIL} across 5 hypotheses. The headline is
\textbf{H2 FAIL}: the D20 aggregate two-cluster structure is
\textbf{NOT robust} at the per-decile granularity.

\begin{table}[h]
\centering
\small
\begin{tabular}{cccccccc}
\hline
decile & $\overline{\mathrm{reward}}$ & best P$_5$ & best P$_6$ & best P$_7$ & best P$_8$ \\
\hline
1  & 0.689 & aero & aero & aero & aero \\
2  & 0.740 & aero & grpo & aero & aero \\
3  & 0.781 & grpo & aero & grpo & grpo \\
4  & 0.814 & aero & grpo & areal & grpo \\
5  & 0.834 & gift & grpo & areal & gift \\
6  & 0.860 & aero & areal & aero & grpo \\
7  & 0.879 & gift & areal & gift & grpo \\
8  & 0.895 & grpo & areal & gift & grpo \\
9  & 0.918 & gift & grpo & areal & gift \\
10 & 0.950 & areal & grpo & areal & grpo \\
\hline
\end{tabular}
\caption{\textbf{Per-decile best method per pillar on N2 same-stack.
\textbf{Best method per pillar VARIES across deciles}; there is no
universally best method. Aero dominates the low-reward regime (deciles
1--2); grpo dominates mid-reward (deciles 3--4, 8); gift dominates
high-mid (5, 7, 9); areal dominates the top decile (10).}}
\label{tab:synth-d21-decile}
\end{table}

\subsubsection{Sharpest findings}\label{sec:synth-d21-findings}

\textbf{F1 (H2 FAIL $\rightarrow$ SHARP).}
The D20 aggregate two-cluster structure \textbf{does NOT survive}
per-decile stratification. Point $\rho$ estimates vary dramatically
($\mathrm{P}_5 \leftrightarrow \mathrm{P}_6$: point $\rho \in [-0.8, +0.8]$
across 10 deciles; $\mathrm{P}_5 \leftrightarrow \mathrm{P}_8$: point
$\rho \in [-0.4, +1.0]$). The two-cluster structure is a
\textbf{pooling artifact} --- it appears at aggregate granularity because
the decile-level structures have heterogeneous signs that average out to
a non-zero aggregate.

\textbf{F2 (H4 PASS).}
Low-reward deciles (1--3) show higher $|\rho|$ than high-reward deciles
(8--10): mean $|\rho|$ is 0.622 (low) vs 0.556 (high). The intuition:
low-reward steps are dominated by failed rollouts where the policy is
unsure, and methods diverge most in those regimes. High-reward steps
are saturated (most methods achieve reward $\sim 1$) and behave
similarly. The discriminative signal lives in the low-reward regime.

\textbf{F3 (cross-decile headline).}
\textbf{No method dominates every pillar in every decile}. The
regime-conditional deployment matrix is:
\begin{itemize}
\item Low-reward (deciles 1--2): \textbf{aero}
\item Mid-reward (deciles 3--4, 8): \textbf{grpo}
\item High-mid (deciles 5, 7, 9): \textbf{gift}
\item Top (decile 10): \textbf{areal}
\end{itemize}
\textbf{Aggregate rankings obscure this regime dependence}; a deployment
that picks one method for the whole reward distribution is
\textbf{operating sub-optimally} on most cells.

\textbf{F4 (methodological).}
With only 4 methods, per-decile $\rho$ confidence intervals are wide
(they always include zero, even when point $\rho = 1.0$). The wide CIs
reflect the structural limitation of having only 4 ranking points, not
absence of signal. The point $\rho$ estimates remain informative; iter-196
surfaces this CI-vs-point discrepancy and recommends that future
cross-pillar audits either pool across more methods (e.g., include
Berkeley variants) or report point estimates with explicit CI caveats.

\subsubsection{Cross-paper coupling}\label{sec:synth-d21-coupling}

\begin{itemize}
\item \textbf{D20 (iter-192)} --- D20's two-cluster structure is
  \textbf{aggregate-only}; iter-196 shows it does not survive
  stratification. The D20 finding should be qualified as holding
  ``at aggregate granularity''.
\item \textbf{D16 (iter-176)} --- D16 measured per-prompt reward stability;
  iter-196 shows best-method-per-pillar varies at the decile level,
  paralleling D16's per-prompt finding that single-method deployment is
  sub-optimal.
\item \textbf{D17 (iter-180)} --- D17 measured paper reproducibility; iter-196's
  wide CIs are themselves a reproducibility concern --- point estimates
  vary but CIs cover everything, suggesting that future cross-pillar
  audits need more than 4 methods to be precise.
\item \textbf{FRONTIER Round 2 (ZVF as signal)} --- D21's regime-conditional
  finding aligns with the frontier synthesis on ZVF as observed signal
  availability rather than difficulty: the regime where methods diverge
  (low reward, low ZVF) is where ZVF-as-signal-starvation matters most.
\end{itemize}

\subsubsection{Operational}\label{sec:synth-d21-operational}
\begin{enumerate}
\item \textbf{QUALIFY the D20 finding}: report that ``two-cluster
  structure'' holds at aggregate granularity but is NOT robust within
  reward-deciles.
\item \textbf{DEPLOY regime-conditionally}: aero for low-reward
  (deciles 1--2), grpo for mid-reward (3--4, 8), gift for high-mid
  (5, 7, 9), areal for top (10).
\item \textbf{REPORT} the per-decile best-method table as
  Table~\ref{tab:synth-d21-decile} in this section.
\item \textbf{WIRE} as CI gate: fails if low-decile $|\rho|$ variance
  drops below 0.55 (the regime-conditional structure collapses).
\item \textbf{EXTEND} in next-iter to per-(decile, task-slice) --- does
  the regime-conditional selection hold across math/coding domains?
\end{enumerate}